\documentclass[11pt,respuestas,a4]{aleph-examen}
% \documentclass[11pt,a5]{aleph-examen}

% -- Paquetes extra
\usepackage{aleph-comandos}
\usepackage{systeme}

% -- Datos
\institucion{Facultad de Ciencias Exactas, Naturales y Ambientales}
\carrera{Catalogo STEM}
\asignatura{Algebra lineal}
\tema{Examen no. 1-2: Sistemas de ecuaciones}
\autor{Andres Merino}
\fecha{Semestre 2026-1}

\logouno[0.16\textwidth]{Logos/logoPUCE_04_ac}
\definecolor{colortext}{HTML}{1957d1}
\definecolor{colordef}{HTML}{1957d1}
\fuente{montserrat}

\begin{document}

\encabezado

\section*{Indicaciones}
\begin{itemize}[leftmargin=*]
\item
	En esta actividad se evalua si el estudiante (Criterio 2.2) aplica diferentes métodos para resolver sistemas de ecuaciones lineales, como el método de eliminación de Gauss y técnicas matriciales.
\item
    Los ejercicios deben ser resueltos a mano, sin ayuda de resúmenes o asistentes computacionales.
\item
	Las resoluciones deben entregarse en hojas de forma física al docente.
\item
	Se encuentra prohibido el uso de cualquier otra fuente de información durante todo el examen.
\item
	En caso de considerar que existe un error en la pregunta o que esta se encuentra mal planteada, se debe indicar cual es el error y justificarlo.
\item
	Todas las soluciones deben estar correctamente redactadas y explicadas.
\item
	Se asignarán 4 puntos de excelencia si, además de tener las soluciones correctas, el examen se destaca por su orden y escritura.
\end{itemize}

\section*{Ejercicios}

\begin{preguntas}

\item 
    Considere el sistema lineal de ecuaciones:
    \[
        \systeme{x+y+2z=2{,}, y+z=1{,}, 2x+4y+6z=6{.}}
    \]
    \begin{enumerate}
        \item Escribe la matriz ampliada del sistema. \puntaje{3}
        \item Aplica reducción de Gauss-Jordan.\puntaje{10}
        \item Analiza la naturaleza de la solución y escríbela.\puntaje{5}
    \end{enumerate}

\begin{respuesta}\hspace{0pt}
\begin{enumerate}
\item La matriz ampliada correspondiente es:
    \[
        \begin{pmatrix}
			 1 & 1 & 2 & | & 2 \\
			 0 & 1 & 1 & | & 1 \\
			 2 & 4 & 6 & | & 6
		\end{pmatrix}.
    \]
\item Aplicando eliminación de Gauss-Jordan:
    \begin{align*}
        \begin{pmatrix}
			1 & 1 & 2 & | & 2 \\
			0 & 1 & 1 & | & 1 \\
			2 & 4 & 6 & | & 6 
		\end{pmatrix}
        &\sim
        \begin{pmatrix}
			1 & 1 & 2 & | & 2 \\
			0 & 1 & 1 & | & 1 \\
			0 & 2 & 2 & | & 2
		\end{pmatrix} && -2F_1 + F_3 \rightarrow F_3 \\
		&\sim
		\begin{pmatrix}
			1 & 1 & 2 & | & 2 \\
			0 & 1 & 1 & | & 1 \\
			0 & 0 & 0 & | & 0
		\end{pmatrix} && -2F_2 + F_3 \rightarrow F_3 \\
		&\sim
		\begin{pmatrix}
			1 & 0 & 1 & | & 1 \\
			0 & 1 & 1 & | & 1 \\
			0 & 0 & 0 & | & 0
		\end{pmatrix} && -F_2 + F_1 \rightarrow F_1.
    \end{align*}
\item Se tiene $\rang(A)=2$ y $\rang(A|b)=2$, por lo tanto el sistema es compatible. Como hay 3 incógnitas y rango 2, el sistema tiene infinitas soluciones (una variable libre).

Si tomamos $z=t$, entonces
\[
    x+z=1\Rightarrow x=1-t,
    \qquad
    y+z=1\Rightarrow y=1-t.
\]
En consecuencia,
\[
    S=\{(x,y,z)=(1-t,1-t,t):\ t\in\R\}.\qedhere
\]
\end{enumerate}
\end{respuesta}

%%%%%%%%%%%%%%%%%%%%%%%%%%%%%%%%%%%%%%%%
%% Inversa
%%%%%%%%%%%%%%%%%%%%%%%%%%%%%%%%%%%%%%%%
\item 
    Considere el sistema lineal de ecuaciones:
    \[
        \systeme{x+y+2 z=2, y+3z=4, 2x-z=1}
    \]
    \begin{enumerate}
        \item Calcule el determinante de la matriz de coeficientes. ¿Esta matriz tiene inversa?\puntaje{7}
        \item Suponga que la inversa de la matriz de coeficientes es:
        \[
            \begin{pmatrix}
               -1 & 1 & 1\\
               6 & -5 & -3\\
               -2 & 2 & 1
            \end{pmatrix},
        \]
        utilícela para calcular la solución del sistema.\puntaje{8}
    \end{enumerate}

\begin{respuesta}\hspace{0pt}
\begin{enumerate}
\item La matriz de coeficientes es:
    \[
        A = \begin{pmatrix}
            1 & 1 & 2 \\
            0 & 1 & 3 \\
            2 & 0 & -1 \\
        \end{pmatrix}.
    \]
    Calculamos el determinante de la matriz de coeficientes:
    \[
        \det(A) = 1\begin{vmatrix} 1 & 3 \\ 0 & -1 \end{vmatrix} - 0\begin{vmatrix} 1 & 2 \\ 0 & -1 \end{vmatrix} + 2\begin{vmatrix} 1 & 2 \\ 1 & 3 \end{vmatrix} = -1 -0 + 2 = 1.
    \]
    Por lo tanto, la matriz de coeficientes tiene inversa.
\item Dado que la matriz de coeficientes tiene inversa, la solución del sistema está dado por $A^{-1}b$, donde $b$ es la matriz de resultados. Por lo tanto:
\[
    \begin{pmatrix}
        x \\ y \\ z
    \end{pmatrix} = \begin{pmatrix}
        -1 & 1 & 1\\
        6 & -5 & -3\\
        -2 & 2 & 1
    \end{pmatrix}\begin{pmatrix}
        2 \\ 4 \\ 1
    \end{pmatrix} = 
    \begin{pmatrix}
        3 \\ -11 \\ 5
    \end{pmatrix}.
\]
Así, la solución del sistema es $x = 3$, $y = -11$ y $z = 5$. \qedhere
\end{enumerate}
\end{respuesta}
%%%%%%%%%%%%%%%%%%%%%%%%%%%%%%%%%%%%%%%%%%
%%%%%%%%%%%%%%%%%%%%%%%%%%%%%%%%%%%%%%%%%%
%%%%%%%%%%%%%%%%%%%%%%%%%%%%%%%%%%%%%%%%%%
\item Sean $r \in \R$. Considere el sistema lineal de ecuaciones:
    \[
        \systeme[xyz]{x + z = 1{,}, 
        y + 2z = 0{,}, 
        x + 2y + rz = 0{,}}
    \]
    Utilizando la eliminación de Gauss-Jordan (indicando lo que realiza paso a paso), determine los valores de $r$ tales que el sistema\puntaje{17}
    \begin{itemize}
        \item tiene una única solución,
        \item tiene infinitas soluciones, y
        \item no tiene solución.
    \end{itemize}

\begin{respuesta}
La matriz ampliada del sistema es:
\[
    \begin{pmatrix}
        1 & 0 & 1 & | & 1 \\
        0 & 1 & 2 & | & 0 \\
        1 & 2 & r & | & 0 \\
    \end{pmatrix}.
\]
Aplicando la eliminación de Gauss-Jordan, obtenemos:
\begin{align*}
    \begin{pmatrix}
        1 & 0 & 1 & | & 1 \\
        0 & 1 & 2 & | & 0 \\
        1 & 2 & r & | & 0 \\
    \end{pmatrix} & \sim  \begin{pmatrix}
        1 & 0 & 1 & | & 1 \\
        0 & 1 & 2 & | & 0 \\
        0 & 2 & r-1 & | & -1 \\
    \end{pmatrix} && - F_1 + F_3 \rightarrow F_3 \\
    & \sim  \begin{pmatrix}
        1 & 0 & 1 & | & 1 \\
        0 & 1 & 2 & | & 0 \\
        0 & 0 & r-5 & | & -1 \\
    \end{pmatrix} && - 2F_2 + F_3 \rightarrow F_3.
\end{align*}
Con esto, tenemos que el rango de la matriz de coeficientes depende del valor de $r$.

Si $r-5=0$ (es decir $r=5$), tenemos que
\[
    \begin{pmatrix}
        1 & 0 & 1 & | & 1 \\
        0 & 1 & 2 & | & 0 \\
        0 & 0 & 0 & | & -1 \\
    \end{pmatrix}
\]
Por lo tanto, $\rang(A)=2$ y $\rang(A|b)=3$, de donde, el sistema es inconsistente, es decir, no tiene solución.

Por otro lado, si $r-5\neq0$ (es decir $r\neq5$), tenemos que $\rang(A)=3$ y $\rang(A|b)=3$, por lo tanto, el sistema es consistente, además, como el sistema tiene 3 incógnitas, se tiene que el sistema posee solución única. 

Finalmente, no existe un valor de $r$ para que el sistema tenga infinitas soluciones.
\end{respuesta}
\end{preguntas}

\end{document}
